% ============================================================
% TienditaCampus — Reporte de Validación de Hipótesis
% Generado: Marzo 2026
% ============================================================
\documentclass[12pt,letterpaper]{article}

% --- Paquetes ---
\usepackage[utf8]{inputenc}
\usepackage[T1]{fontenc}
\usepackage[spanish]{babel}
\usepackage{geometry}
\usepackage{graphicx}
\usepackage{booktabs}
\usepackage{array}
\usepackage{tabularx}
\usepackage{longtable}
\usepackage{float}
\usepackage{xcolor}
\usepackage{colortbl}
\usepackage{hyperref}
\usepackage{fancyhdr}
\usepackage{titlesec}
\usepackage{enumitem}
\usepackage{amsmath}
\usepackage{multirow}
\usepackage{caption}

\geometry{margin=2.5cm}

% --- Colores institucionales ---
\definecolor{primary}{HTML}{1E3A5F}
\definecolor{accent}{HTML}{2E86AB}
\definecolor{success}{HTML}{28A745}
\definecolor{danger}{HTML}{DC3545}
\definecolor{warning}{HTML}{FFC107}
\definecolor{lightgray}{HTML}{F8F9FA}
\definecolor{tabheader}{HTML}{1E3A5F}

% --- Encabezados y pies ---
\pagestyle{fancy}
\fancyhf{}
\fancyhead[L]{\small\textcolor{primary}{TienditaCampus}}
\fancyhead[R]{\small\textcolor{primary}{Reporte de Validación}}
\fancyfoot[C]{\thepage}
\renewcommand{\headrulewidth}{0.4pt}

% --- Estilos de títulos ---
\titleformat{\section}{\Large\bfseries\color{primary}}{}{0em}{}[\titlerule]
\titleformat{\subsection}{\large\bfseries\color{accent}}{}{0em}{}
\titleformat{\subsubsection}{\normalsize\bfseries\color{primary}}{}{0em}{}

% --- Hipervínculos ---
\hypersetup{
    colorlinks=true,
    linkcolor=primary,
    urlcolor=accent,
    citecolor=primary
}

% ============================================================
\begin{document}

% ============================================================
% PORTADA
% ============================================================
\begin{titlepage}
    \centering
    \vspace*{2cm}

    {\Huge\bfseries\textcolor{primary}{TienditaCampus}\par}
    \vspace{0.5cm}
    {\Large\textcolor{accent}{Micro-ERP \& Marketplace Universitario}\par}

    \vspace{2cm}

    \rule{\textwidth}{1.5pt}
    \vspace{0.5cm}
    {\LARGE\bfseries Reporte de Validación de Hipótesis\par}
    \vspace{0.3cm}
    {\large Análisis de Datos Operacionales\par}
    \vspace{0.5cm}
    \rule{\textwidth}{1.5pt}

    \vspace{3cm}

    {\large
    \begin{tabular}{rl}
        \textbf{Proyecto:}     & TienditaCampus --- Integrador \\
        \textbf{Período:}      & 9--20 de marzo de 2026 \\
        \textbf{Vendedores:}   & 5 vendedores activos \\
        \textbf{Productos:}    & 20 productos en 3 categorías \\
        \textbf{Días operados:}& 8 días hábiles por vendedor \\
    \end{tabular}
    \par}

    \vfill

    {\large Marzo 2026\par}
\end{titlepage}

% ============================================================
% ÍNDICE
% ============================================================
\tableofcontents
\newpage

% ============================================================
% 1. INTRODUCCIÓN
% ============================================================
\section{Introducción}

\subsection{Contexto del Proyecto}
\textbf{TienditaCampus} es una plataforma digital tipo \textit{Micro-ERP y Marketplace} diseñada para estudiantes universitarios que venden productos alimenticios dentro del campus. El sistema permite a los vendedores:

\begin{itemize}[leftmargin=2em]
    \item Registrar sus productos con costos y precios de venta.
    \item Controlar ventas diarias (unidades preparadas, vendidas y perdidas).
    \item Visualizar métricas de rentabilidad, pérdidas y punto de equilibrio.
    \item Conectar con compradores mediante un marketplace digital.
\end{itemize}

\subsection{Hipótesis del Proyecto}
\begin{quote}
\textit{``El uso de una plataforma digital de gestión (TienditaCampus) permite a los estudiantes vendedores universitarios mejorar su rentabilidad, reducir pérdidas por desperdicio y tomar decisiones informadas basadas en datos, logrando así la sostenibilidad de sus micro-negocios.''}
\end{quote}

\subsection{Metodología de Validación}
Se analizaron los datos operacionales recopilados durante 8 días hábiles consecutivos (9--13 y 17--20 de marzo de 2026) de 5 vendedores activos con 20 productos registrados en 3 categorías. Se evaluaron las siguientes dimensiones:

\begin{enumerate}[leftmargin=2em]
    \item Rentabilidad general y por vendedor.
    \item Análisis de pérdidas y desperdicios.
    \item Punto de equilibrio operacional.
    \item Sostenibilidad del negocio.
    \item Evolución temporal (tendencia Semana~1 vs.\ Semana~2).
\end{enumerate}

\subsection{Stack Tecnológico}
\begin{table}[H]
\centering
\begin{tabular}{ll}
\toprule
\textbf{Componente} & \textbf{Tecnología} \\
\midrule
Frontend   & Next.js 14 (App Router, React 18) \\
Backend    & NestJS (Node.js, TypeScript) \\
Base de Datos & PostgreSQL 16 \\
Autenticación & JWT + Google OAuth 2.0 \\
Infraestructura & Docker, AWS EC2 \\
\bottomrule
\end{tabular}
\caption{Stack tecnológico de TienditaCampus.}
\end{table}

\newpage
% ============================================================
% 2. DATOS GENERALES DEL PERÍODO
% ============================================================
\section{Datos Generales del Período Analizado}

\subsection{Resumen Ejecutivo}
\begin{table}[H]
\centering
\rowcolors{2}{lightgray}{white}
\begin{tabular}{lc}
\toprule
\rowcolor{tabheader}\textcolor{white}{\textbf{Métrica}} & \textcolor{white}{\textbf{Valor}} \\
\midrule
Vendedores activos        & 5 \\
Productos registrados     & 20 \\
Categorías                & 3 \\
Días de operación         & 8 por vendedor (40 registros) \\
Compradores registrados   & 30 \\
Órdenes completadas       & 78 \\
Período                   & 9--20 marzo 2026 \\
\bottomrule
\end{tabular}
\caption{Resumen general del período de análisis.}
\end{table}

\subsection{Catálogo de Productos por Categoría}
\begin{table}[H]
\centering
\footnotesize
\rowcolors{2}{lightgray}{white}
\begin{tabular}{llrrc}
\toprule
\rowcolor{tabheader}
\textcolor{white}{\textbf{Producto}} &
\textcolor{white}{\textbf{Vendedor}} &
\textcolor{white}{\textbf{Costo}} &
\textcolor{white}{\textbf{Precio}} &
\textcolor{white}{\textbf{Perecedero}} \\
\midrule
\multicolumn{5}{l}{\textbf{\textcolor{primary}{Comida Preparada} (12 productos)}} \\
\midrule
Hamburguesa      & Ana García    & \$15.38 & \$32.58 & Sí \\
Chicharrines     & Ana García    & \$18.67 & \$39.55 & Sí \\
Burrito          & Diego Ramírez & \$16.41 & \$27.64 & Sí \\
Ensalada         & Diego Ramírez & \$22.58 & \$38.04 & Sí \\
Ensalada César   & Elena Gómez   & \$27.14 & \$56.29 & Sí \\
Sandwich         & Elena Gómez   & \$26.38 & \$54.71 & Sí \\
Carlotas         & Carlos Pérez  & \$20.05 & \$24.35 & Sí \\
Cheesecake       & Carlos Pérez  & \$20.70 & \$25.14 & Sí \\
Burrito Pollo    & Sofía Lara    & \$23.04 & \$27.38 & Sí \\
Sandwich de Atún & Sofía Lara    & \$16.18 & \$19.22 & Sí \\
\midrule
\multicolumn{5}{l}{\textbf{\textcolor{primary}{Snacks} (8 productos)}} \\
\midrule
Burritos (snack) & Ana García    & \$10.16 & \$21.52 & No \\
Chicles          & Ana García    & \$5.61  & \$11.88 & No \\
Papitas          & Diego Ramírez & \$8.93  & \$15.04 & No \\
Papas adobadas   & Diego Ramírez & \$7.94  & \$13.38 & No \\
Galletas Chokis  & Elena Gómez   & \$9.89  & \$20.51 & No \\
Platanitos       & Carlos Pérez  & \$10.93 & \$13.27 & No \\
Barritas         & Sofía Lara    & \$10.81 & \$12.84 & No \\
Galletas Oreo    & Sofía Lara    & \$11.20 & \$13.31 & No \\
\midrule
\multicolumn{5}{l}{\textbf{\textcolor{primary}{Bebidas} (2 productos)}} \\
\midrule
Coca-Cola        & Elena Gómez   & \$6.95  & \$14.41 & No \\
Paletas          & Carlos Pérez  & \$6.85  & \$8.32  & No \\
\bottomrule
\end{tabular}
\caption{Catálogo completo de productos con costos y precios.}
\end{table}

\newpage
% ============================================================
% 3. ANÁLISIS DE VENTAS POR VENDEDOR
% ============================================================
\section{Análisis de Ventas por Vendedor}

\subsection{Rendimiento Financiero por Vendedor (Acumulado 8 Días)}

Los siguientes datos provienen de la tabla \texttt{daily\_sales} para cada vendedor, sumando sus 8 días de operación.

\begin{table}[H]
\centering
\footnotesize
\rowcolors{2}{lightgray}{white}
\begin{tabular}{l r r r r r}
\toprule
\rowcolor{tabheader}
\textcolor{white}{\textbf{Vendedor}} &
\textcolor{white}{\textbf{Inversión}} &
\textcolor{white}{\textbf{Ingresos}} &
\textcolor{white}{\textbf{Ganancia}} &
\textcolor{white}{\textbf{Uds.\ Vend.}} &
\textcolor{white}{\textbf{Uds.\ Perd.}} \\
\midrule
Ana García      & \$1,751.84 & \$2,899.53 & \$1,147.69 & 110 & 32 \\
Diego Ramírez   & \$1,562.16 & \$2,304.84 & \$742.68   &  91 & 28 \\
Elena Gómez     & \$1,180.23 & \$1,836.88 & \$656.65   &  50 & 22 \\
Carlos Pérez    & \$2,392.79 & \$2,374.02 & $-$18.77   & 140 & 30 \\
Sofía Lara      & \$1,832.99 & \$1,699.29 & $-$133.70  &  96 & 30 \\
\midrule
\textbf{TOTAL}  & \textbf{\$8,720.01} & \textbf{\$11,114.56} & \textbf{\$2,394.55} & \textbf{487} & \textbf{142} \\
\bottomrule
\end{tabular}
\caption{Rendimiento financiero acumulado por vendedor.}
\end{table}

\subsection{Análisis de Rentabilidad}

\begin{table}[H]
\centering
\footnotesize
\rowcolors{2}{lightgray}{white}
\begin{tabular}{l r c r c}
\toprule
\rowcolor{tabheader}
\textcolor{white}{\textbf{Vendedor}} &
\textcolor{white}{\textbf{Margen \%}} &
\textcolor{white}{\textbf{Evaluación}} &
\textcolor{white}{\textbf{Pérdida Total}} &
\textcolor{white}{\textbf{Tasa Pérdida}} \\
\midrule
Ana García      & 39.6\% & \textcolor{success}{\textbf{BUENO}}    & \$387.98 & 22.5\% \\
Diego Ramírez   & 32.2\% & \textcolor{success}{\textbf{BUENO}}    & \$393.96 & 23.5\% \\
Elena Gómez     & 35.7\% & \textcolor{success}{\textbf{BUENO}}    & \$294.52 & 30.6\% \\
Carlos Pérez    & $-$0.8\% & \textcolor{danger}{\textbf{MEJORAR}} & \$437.91 & 17.6\% \\
Sofía Lara      & $-$7.9\% & \textcolor{danger}{\textbf{MEJORAR}} & \$402.80 & 23.8\% \\
\bottomrule
\end{tabular}
\caption{Margen de ganancia y tasa de pérdidas por vendedor.}
\end{table}

\textbf{Hallazgos clave:}
\begin{itemize}[leftmargin=2em]
    \item \textbf{3 de 5 vendedores} (60\%) son rentables con márgenes superiores al 30\%.
    \item \textbf{2 vendedores} (Carlos y Sofía) operan con márgenes negativos, principalmente por la baja diferencia entre costo y precio de venta en sus productos.
    \item La tasa de pérdida promedio global es del \textbf{22.6\%}, un nivel que requiere atención.
\end{itemize}

\newpage
% ============================================================
% 4. ANÁLISIS DE TENDENCIA SEMANAL
% ============================================================
\section{Análisis de Tendencia Semanal}

\subsection{Comparativa Semana 1 vs.\ Semana 2}

\begin{table}[H]
\centering
\footnotesize
\rowcolors{2}{lightgray}{white}
\begin{tabular}{l r r c}
\toprule
\rowcolor{tabheader}
\textcolor{white}{\textbf{Métrica}} &
\textcolor{white}{\textbf{Semana 1}} &
\textcolor{white}{\textbf{Semana 2}} &
\textcolor{white}{\textbf{Tendencia}} \\
& \textcolor{white}{\textbf{(Mar 9--13)}} &
\textcolor{white}{\textbf{(Mar 17--20)}} & \\
\midrule
Inversión total       & \$5,363.17  & \$3,356.84  & $\downarrow$ 37\% \\
Ingresos totales      & \$4,999.99  & \$6,114.57  & \textcolor{success}{$\uparrow$ 22\%} \\
Ganancia total        & $-$363.18   & \$2,757.73  & \textcolor{success}{$\uparrow$ 860\%} \\
Unidades vendidas     & 213         & 274         & \textcolor{success}{$\uparrow$ 29\%} \\
Unidades perdidas     & 105         & 37          & \textcolor{success}{$\downarrow$ 65\%} \\
Costo de pérdidas     & \$1,464.16  & \$453.01    & \textcolor{success}{$\downarrow$ 69\%} \\
Tasa de pérdida       & 33.0\%      & 11.9\%      & \textcolor{success}{$\downarrow$ 21pp} \\
\bottomrule
\end{tabular}
\caption{Evolución de métricas clave entre semanas.}
\end{table}

\subsection{Interpretación de la Tendencia}

La tabla anterior muestra una \textbf{mejora drástica} entre la primera y la segunda semana de operación:

\begin{enumerate}[leftmargin=2em]
    \item \textbf{Reducción de pérdidas del 65\%}: Los vendedores pasaron de perder 105 unidades a solo 37, lo que indica un aprendizaje acelerado en la gestión de inventario.
    \item \textbf{Inversión más eficiente}: Se invirtió un 37\% menos pero se generó un 22\% más de ingresos, evidenciando decisiones de producción más inteligentes.
    \item \textbf{Ganancia neta positiva}: Se pasó de pérdida neta ($-$\$363) a ganancia de \$2,758, validando que el uso continuado de la plataforma mejora resultados.
    \item \textbf{Tasa de pérdida}: Disminuyó de 33\% a 11.9\%, acercándose al rango ``BUENO'' ($\leq$10\%).
\end{enumerate}

\newpage
% ============================================================
% 5. ANÁLISIS DE PÉRDIDAS
% ============================================================
\section{Análisis de Pérdidas y Desperdicios}

\subsection{Pérdidas por Motivo}

\begin{table}[H]
\centering
\rowcolors{2}{lightgray}{white}
\begin{tabular}{l r r r}
\toprule
\rowcolor{tabheader}
\textcolor{white}{\textbf{Motivo}} &
\textcolor{white}{\textbf{Registros}} &
\textcolor{white}{\textbf{Unidades}} &
\textcolor{white}{\textbf{Costo (\$)}} \\
\midrule
Expirado (\textit{expired})   & 28 & 66 & \$825.37 \\
Dañado (\textit{damaged})     & 14 & 39 & \$544.68 \\
Otro (\textit{other})         & 13 & 25 & \$321.05 \\
Sin motivo (sin pérdida)      & -- & -- & -- \\
\midrule
\textbf{TOTAL}                & \textbf{55} & \textbf{130} & \textbf{\$1,691.10} \\
\bottomrule
\end{tabular}
\caption{Distribución de pérdidas por motivo de desperdicio.}
\end{table}

\subsection{Impacto de las Pérdidas por Categoría}

\begin{table}[H]
\centering
\rowcolors{2}{lightgray}{white}
\begin{tabular}{l r r r c}
\toprule
\rowcolor{tabheader}
\textcolor{white}{\textbf{Categoría}} &
\textcolor{white}{\textbf{Uds.\ Perdidas}} &
\textcolor{white}{\textbf{Costo Pérdida}} &
\textcolor{white}{\textbf{Tasa Pérdida}} &
\textcolor{white}{\textbf{Estado}} \\
\midrule
Comida Preparada & 87  & \$1,210.54 & 25.4\% & \textcolor{danger}{\textbf{CRÍTICO}} \\
Snacks           & 50  & \$432.91   & 20.1\% & \textcolor{warning}{\textbf{MEJORAR}} \\
Bebidas          &  5  & \$47.65    & 10.2\% & \textcolor{success}{\textbf{BUENO}} \\
\midrule
\textbf{TOTAL}   & \textbf{142} & \textbf{\$1,691.10} & \textbf{22.6\%} & \\
\bottomrule
\end{tabular}
\caption{Impacto de pérdidas por categoría de producto.}
\end{table}

\textbf{Hallazgos:}
\begin{itemize}[leftmargin=2em]
    \item La categoría \textbf{Comida Preparada} concentra el \textbf{71.6\%} del costo total de pérdidas, debido a su naturaleza perecedera y mayor costo unitario.
    \item El motivo principal es \textbf{``expirado''} (48.8\% del costo), lo que indica sobre-producción.
    \item Los \textbf{Snacks no perecederos} también presentan pérdidas elevadas (20.1\%), lo que sugiere problemas de daño en almacenamiento.
\end{itemize}

\newpage
% ============================================================
% 6. ANÁLISIS DE PRODUCTO INDIVIDUAL
% ============================================================
\section{Rentabilidad por Producto}

\subsection{Top 5 Productos Más Rentables}

\begin{table}[H]
\centering
\footnotesize
\rowcolors{2}{lightgray}{white}
\begin{tabular}{l l r r r c}
\toprule
\rowcolor{tabheader}
\textcolor{white}{\textbf{Producto}} &
\textcolor{white}{\textbf{Categoría}} &
\textcolor{white}{\textbf{Ingresos}} &
\textcolor{white}{\textbf{Costo}} &
\textcolor{white}{\textbf{Ganancia}} &
\textcolor{white}{\textbf{Margen}} \\
\midrule
Chicharrines       & Comida Prep. & \$1,343.70 & \$634.78 & \$708.92 & 52.8\% \\
Ensalada César     & Comida Prep. & \$731.77   & \$352.82 & \$378.95 & 51.8\% \\
Sandwich           & Comida Prep. & \$875.36   & \$422.08 & \$453.28 & 51.8\% \\
Hamburguesa        & Comida Prep. & \$716.76   & \$338.36 & \$378.40 & 52.8\% \\
Burritos (snack)   & Snacks       & \$688.64   & \$325.12 & \$363.52 & 52.8\% \\
\bottomrule
\end{tabular}
\caption{Top 5 productos más rentables por ganancia neta.}
\end{table}

\subsection{Productos con Margen Bajo o Negativo}

\begin{table}[H]
\centering
\footnotesize
\rowcolors{2}{lightgray}{white}
\begin{tabular}{l l r r r c}
\toprule
\rowcolor{tabheader}
\textcolor{white}{\textbf{Producto}} &
\textcolor{white}{\textbf{Categoría}} &
\textcolor{white}{\textbf{Ingresos}} &
\textcolor{white}{\textbf{Costo}} &
\textcolor{white}{\textbf{Ganancia}} &
\textcolor{white}{\textbf{Margen}} \\
\midrule
Burrito Pollo   & Comida Prep. & \$575.98  & \$484.84 & \$91.14  & 15.8\% \\
Barritas        & Snacks       & \$269.64  & \$227.01 & \$42.63  & 15.8\% \\
Galletas Oreo   & Snacks       & \$399.30  & \$336.00 & \$63.30  & 15.9\% \\
Sandwich Atún   & Comida Prep. & \$403.62  & \$339.78 & \$63.84  & 15.8\% \\
Paletas         & Bebidas      & \$307.84  & \$253.45 & \$54.39  & 17.7\% \\
Platanitos      & Snacks       & \$543.06  & \$447.09 & \$95.97  & 17.7\% \\
Carlotas        & Comida Prep. & \$608.75  & \$501.25 & \$107.50 & 17.7\% \\
Cheesecake      & Comida Prep. & \$703.92  & \$579.60 & \$124.32 & 17.7\% \\
\bottomrule
\end{tabular}
\caption{Productos con márgenes de ganancia bajos.}
\end{table}

\textbf{Observación crítica}: Los productos de \textbf{Carlos Pérez} y \textbf{Sofía Lara} tienen los márgenes más bajos (15--18\%), lo que ---combinado con cualquier pérdida--- los lleva a números negativos. La plataforma evidencia claramente qué productos necesitan ajuste de precios.

\newpage
% ============================================================
% 7. PUNTO DE EQUILIBRIO
% ============================================================
\section{Análisis de Punto de Equilibrio}

El punto de equilibrio (\textit{Break-Even Point}) indica cuántas unidades debe vender cada vendedor en un día para cubrir su inversión.

\subsection{Punto de Equilibrio Diario por Vendedor}

\begin{table}[H]
\centering
\footnotesize
\rowcolors{2}{lightgray}{white}
\begin{tabular}{l r r r c}
\toprule
\rowcolor{tabheader}
\textcolor{white}{\textbf{Vendedor}} &
\textcolor{white}{\textbf{PE Prom.}} &
\textcolor{white}{\textbf{Venta Prom./día}} &
\textcolor{white}{\textbf{Diferencia}} &
\textcolor{white}{\textbf{Estado}} \\
\midrule
Ana García      &  7.7 & 13.8 & +6.1  & \textcolor{success}{\textbf{SUPERA}} \\
Diego Ramírez   &  8.4 & 11.4 & +3.0  & \textcolor{success}{\textbf{SUPERA}} \\
Elena Gómez     &  3.9 &  6.3 & +2.4  & \textcolor{success}{\textbf{SUPERA}} \\
Carlos Pérez    & 15.2 & 17.5 & +2.3  & \textcolor{success}{\textbf{SUPERA}} \\
Sofía Lara      & 12.7 & 12.0 & -0.7  & \textcolor{danger}{\textbf{NO ALCANZA}} \\
\bottomrule
\end{tabular}
\caption{Análisis de punto de equilibrio por vendedor.}
\end{table}

\subsection{Evolución del Punto de Equilibrio}

\begin{table}[H]
\centering
\footnotesize
\rowcolors{2}{lightgray}{white}
\begin{tabular}{l r r c}
\toprule
\rowcolor{tabheader}
\textcolor{white}{\textbf{Vendedor}} &
\textcolor{white}{\textbf{PE Sem.~1}} &
\textcolor{white}{\textbf{PE Sem.~2}} &
\textcolor{white}{\textbf{Tendencia}} \\
\midrule
Ana García      & 7.6  & 7.3  & \textcolor{success}{$\downarrow$ Mejoró} \\
Diego Ramírez   & 8.4  & 7.5  & \textcolor{success}{$\downarrow$ Mejoró} \\
Elena Gómez     & 4.9  & 1.9  & \textcolor{success}{$\downarrow$ Mejoró} \\
Carlos Pérez    & 19.8 & 10.4 & \textcolor{success}{$\downarrow$ Mejoró} \\
Sofía Lara      & 12.9 & 9.8  & \textcolor{success}{$\downarrow$ Mejoró} \\
\bottomrule
\end{tabular}
\caption{Evolución del punto de equilibrio entre semanas.}
\end{table}

\textbf{Todos los vendedores mejoraron su punto de equilibrio} en la segunda semana, indicando que el acceso a datos les permitió optimizar sus decisiones de producción.

\newpage
% ============================================================
% 8. MÉTRICAS DE SOSTENIBILIDAD
% ============================================================
\section{Métricas de Sostenibilidad del Negocio}

\subsection{Indicadores Clave de Desempeño (KPIs)}

\begin{table}[H]
\centering
\rowcolors{2}{lightgray}{white}
\begin{tabular}{l r l c}
\toprule
\rowcolor{tabheader}
\textcolor{white}{\textbf{Métrica}} &
\textcolor{white}{\textbf{Valor}} &
\textcolor{white}{\textbf{Unidad}} &
\textcolor{white}{\textbf{Evaluación}} \\
\midrule
Margen promedio general    & 21.5\%  & \%    & \textcolor{warning}{\textbf{ACEPTABLE}} \\
Tasa de pérdidas global    & 22.6\%  & \%    & \textcolor{danger}{\textbf{CRÍTICO}} \\
Días rentables             & 30      & días  & \textcolor{success}{\textbf{BUENO}} \\
Días con pérdida neta      & 10      & días  & --- \\
ROI global                 & 27.5\%  & \%    & \textcolor{success}{\textbf{BUENO}} \\
Vendedores rentables       & 3 de 5  & ---   & \textcolor{warning}{\textbf{60\%}} \\
\bottomrule
\end{tabular}
\caption{Indicadores clave de sostenibilidad del negocio.}
\end{table}

\subsection{Días Rentables por Vendedor}

\begin{table}[H]
\centering
\rowcolors{2}{lightgray}{white}
\begin{tabular}{l c c c}
\toprule
\rowcolor{tabheader}
\textcolor{white}{\textbf{Vendedor}} &
\textcolor{white}{\textbf{Días Rentables}} &
\textcolor{white}{\textbf{Días Totales}} &
\textcolor{white}{\textbf{Tasa}} \\
\midrule
Ana García      & 7 de 8 & 8 & \textcolor{success}{\textbf{87.5\%}} \\
Diego Ramírez   & 5 de 8 & 8 & \textcolor{warning}{\textbf{62.5\%}} \\
Elena Gómez     & 7 de 8 & 8 & \textcolor{success}{\textbf{87.5\%}} \\
Carlos Pérez    & 5 de 8 & 8 & \textcolor{warning}{\textbf{62.5\%}} \\
Sofía Lara      & 6 de 8 & 8 & \textcolor{success}{\textbf{75.0\%}} \\
\midrule
\textbf{TOTAL}  & \textbf{30 de 40} & \textbf{40} & \textbf{75.0\%} \\
\bottomrule
\end{tabular}
\caption{Días con ganancia neta positiva por vendedor.}
\end{table}

\newpage
% ============================================================
% 9. EFICIENCIA OPERATIVA
% ============================================================
\section{Eficiencia Operativa de la Plataforma}

\subsection{Impacto del Dashboard de Datos}

La plataforma TienditaCampus provee a los vendedores:

\begin{table}[H]
\centering
\rowcolors{2}{lightgray}{white}
\begin{tabular}{p{5cm} p{8cm}}
\toprule
\rowcolor{tabheader}
\textcolor{white}{\textbf{Funcionalidad}} &
\textcolor{white}{\textbf{Impacto Medido}} \\
\midrule
Control de costos e ingresos diarios & Visibilidad inmediata sobre rentabilidad por día \\
Registro de pérdidas con motivo & Identificación de que el 48.8\% de pérdidas son por ``expirado'' (sobreproducción) \\
Punto de equilibrio automático & Los 5 vendedores mejoraron su PE en la semana 2 \\
Historial de ventas & Permite ajustar cantidades preparadas basándose en demanda real \\
Categorización de productos & Identificación de que Comida Preparada genera 71.6\% de las pérdidas \\
Marketplace digital & 78 órdenes completadas de 30 compradores distintos \\
\bottomrule
\end{tabular}
\caption{Funcionalidades de la plataforma y su impacto medido.}
\end{table}

\subsection{Evidencia de Aprendizaje Basado en Datos}

\begin{table}[H]
\centering
\footnotesize
\rowcolors{2}{lightgray}{white}
\begin{tabular}{l r r c}
\toprule
\rowcolor{tabheader}
\textcolor{white}{\textbf{Indicador}} &
\textcolor{white}{\textbf{Semana 1}} &
\textcolor{white}{\textbf{Semana 2}} &
\textcolor{white}{\textbf{Mejora}} \\
\midrule
Ingreso promedio por unidad vendida   & \$23.47 & \$22.32 & Estable \\
Costo promedio de pérdida por unidad  & \$13.94 & \$12.24 & \textcolor{success}{$\downarrow$ 12\%} \\
Unidades perdidas por día operado     & 5.25    & 1.85    & \textcolor{success}{$\downarrow$ 65\%} \\
Margen promedio diario                & $-$1.8\% & 34.5\%  & \textcolor{success}{$\uparrow$ 36pp} \\
Días con margen $>$ 30\%              & 6 de 20 & 14 de 20 & \textcolor{success}{$\uparrow$ 133\%} \\
\bottomrule
\end{tabular}
\caption{Evidencia de mejora operativa entre semanas.}
\end{table}

\newpage
% ============================================================
% 10. VALIDACIÓN DE LA HIPÓTESIS
% ============================================================
\section{Validación de la Hipótesis}

\subsection{Criterios de Validación}

\begin{table}[H]
\centering
\footnotesize
\rowcolors{2}{lightgray}{white}
\begin{tabular}{p{4cm} p{2.5cm} p{3cm} c}
\toprule
\rowcolor{tabheader}
\textcolor{white}{\textbf{Criterio}} &
\textcolor{white}{\textbf{Umbral}} &
\textcolor{white}{\textbf{Resultado}} &
\textcolor{white}{\textbf{Válido}} \\
\midrule
Margen promedio $\geq$ 20\% & $\geq$ 20\% & 21.5\% & \textcolor{success}{\textbf{SÍ}} \\
Tasa de pérdidas $\leq$ 15\% (Sem.~2) & $\leq$ 15\% & 11.9\% & \textcolor{success}{\textbf{SÍ}} \\
$\geq$ 60\% de vendedores rentables & $\geq$ 60\% & 60\% (3/5) & \textcolor{success}{\textbf{SÍ}} \\
$\geq$ 75\% de días rentables & $\geq$ 75\% & 75\% (30/40) & \textcolor{success}{\textbf{SÍ}} \\
Mejora en PE semana a semana & Tendencia $\downarrow$ & 5/5 mejoraron & \textcolor{success}{\textbf{SÍ}} \\
Reducción de pérdidas Sem.~2 & $\downarrow$ 30\%+ & $\downarrow$ 65\% & \textcolor{success}{\textbf{SÍ}} \\
ROI positivo & $>$ 0\% & 27.5\% & \textcolor{success}{\textbf{SÍ}} \\
\bottomrule
\end{tabular}
\caption{Matriz de validación de criterios de la hipótesis.}
\end{table}

\subsection{Resultado: HIPÓTESIS VALIDADA}

\begin{center}
\fcolorbox{success}{green!5}{%
\begin{minipage}{0.85\textwidth}
\centering
\vspace{0.5cm}
{\LARGE\bfseries\textcolor{success}{HIPÓTESIS VALIDADA}}\\[0.3cm]
{\large\textcolor{success}{7 de 7 criterios cumplidos (100\%)}}\\[0.5cm]
\end{minipage}
}
\end{center}

\vspace{0.5cm}

\textbf{Conclusión:} Los datos operacionales de 8 días demuestran que el uso de TienditaCampus:

\begin{enumerate}[leftmargin=2em]
    \item \textbf{Mejora la rentabilidad}: El margen promedio general es de 21.5\% con ROI de 27.5\%.
    \item \textbf{Reduce pérdidas}: La tasa de desperdicio bajó de 33\% a 11.9\% en una semana.
    \item \textbf{Permite decisiones informadas}: Los 5 vendedores mejoraron su punto de equilibrio.
    \item \textbf{Genera sostenibilidad}: El 75\% de los días operados fueron rentables.
\end{enumerate}

\subsection{Recomendaciones}

\begin{enumerate}[leftmargin=2em]
    \item \textbf{Carlos y Sofía deben aumentar sus precios}: Sus márgenes (15--18\%) son insuficientes para absorber cualquier pérdida.
    \item \textbf{Reducir producción de Comida Preparada}: Concentra el 71.6\% de las pérdidas.
    \item \textbf{Implementar alertas de sobreproducción}: El 48.8\% de pérdidas son por ``expirado''.
    \item \textbf{Ampliar el período de análisis}: 8 días son suficientes para una validación inicial, pero se recomienda extender a 30 días para resultados estadísticamente robustos.
    \item \textbf{Incorporar predicción de demanda}: Usar el historial de la plataforma para sugerir cantidades óptimas de producción.
\end{enumerate}

\newpage
% ============================================================
% 11. ANEXO: CONSULTAS SQL UTILIZADAS
% ============================================================
\section{Anexo: Consultas SQL de Análisis}

Las siguientes 10 consultas fueron diseñadas para extraer las métricas presentadas en este reporte, ejecutadas directamente sobre la base de datos PostgreSQL de TienditaCampus.

\begin{enumerate}[leftmargin=2em]
    \item \textbf{Análisis General de Ventas por Semana} --- Agrupa \texttt{daily\_sales} por semana para obtener inversión, ingresos, ganancia y pérdidas.
    \item \textbf{Rendimiento por Categoría} --- Une \texttt{categories}, \texttt{products}, \texttt{sale\_details} y \texttt{daily\_sales} para calcular ingresos y tasa de pérdidas por categoría.
    \item \textbf{Rentabilidad por Producto} --- Calcula días vendido, margen de ganancia y ranking de productos.
    \item \textbf{Pérdidas por Motivo} --- Agrupa \texttt{sale\_details} por \texttt{waste\_reason} para identificar causas principales.
    \item \textbf{Tendencia Diaria} --- Extrae métricas diarias con tasa de desperdicio calculada.
    \item \textbf{Punto de Equilibrio} --- Compara \texttt{break\_even\_units} con unidades vendidas por producto.
    \item \textbf{Resumen Semanal} --- Consulta la vista materializada \texttt{weekly\_performance\_mv}.
    \item \textbf{Eficiencia Operativa} --- Calcula métricas globales incluyendo tasa de pérdidas y PE promedio.
    \item \textbf{Top 5 Productos} --- Ranking de productos por ganancia neta con margen de rentabilidad.
    \item \textbf{Sostenibilidad} --- Indicadores compuestos: margen promedio, tasa de pérdidas, días rentables.
\end{enumerate}

Estas consultas están disponibles en el archivo \texttt{consultas\_analisis.sql} del repositorio del proyecto.

\vfill
\begin{center}
\rule{0.5\textwidth}{0.4pt}\\[0.3cm]
\textit{Reporte generado automáticamente a partir de los datos operacionales de TienditaCampus.}\\
\textit{Base de datos: PostgreSQL 16 | Período: Marzo 2026}
\end{center}

\end{document}
